%
% Appendix B - Selected Code Listings
%
\chapter{Selected Code Listings} \label{app:code-listings}

This appendix gives the full, unabridged source of a few backend and frontend units that Chapters~\ref{cap:backend} and~\ref{cap:frontend} show only in an abbreviated form, to keep those chapters within the report's page budget.

\section{The Weekend Rule Definition} \label{app:weekend-rule}

Referenced from Section~\ref{sec:rule-engine}.

\begin{lstlisting}[caption={[The Weekend rule definition, in full.]{The Weekend rule definition, in full (\texttt{server/calendar/rules/definitions/weekend.rule.ts}).}}, label={lst:weekend-rule-full}]
export const weekendRule: IRuleDefinition = {
  id: "weekend",
  fields: [
    { id: "checkStart", type: "checkbox", defaultValue: true },
    { id: "checkEnd",   type: "checkbox", defaultValue: true },
    { id: "checkRange", type: "checkbox", defaultValue: false },
  ],
  translations: {
    pt: {
      label: "Fim de semana",
      fields: { checkStart: "Verificar início", checkEnd: "Verificar fim", checkRange: "Verificar período" },
      fieldLabels: { start: "Início", end: "Fim", range: "Período" },
      messages: {
        startIsWeekend: "O início coincide com um fim de semana",
        endIsWeekend: "O fim coincide com um fim de semana",
        rangeHasWeekend: "O período inclui um dia de fim de semana",
      },
    },
  },
  validate(_event, occurrence, config) {
    return runDayRangeCheck<true>({
      occurrence,
      config: config as WeekendConfig,
      matchDay: (day) => (isWeekend(day) ? true : undefined),
      toViolation: (field, date) => ({ fieldLabelKey: field, date, messageKey: MESSAGE_KEYS[field] }),
    });
  },
};
\end{lstlisting}

\clearpage
\section{The Global Rule-Validation Dispatcher} \label{app:collect-issues}

Referenced from Section~\ref{sec:rule-engine}.

\begin{lstlisting}[caption={[The global rule-validation dispatcher, in full.]{The global rule-validation dispatcher, in full (\texttt{server/calendar/services/calendar.service.ts}).}}, label={lst:collect-issues-full}]
private collectEventRuleIssues(
  event: IEvent,
  academicYearStart: number,
  allEvents: IEvent[],
  vacations: IVacationPeriod[] = [],
): ICalendarRuleIssue[] {
  const context = { allEvents, academicYearStart, vacations };
  const issues: ICalendarRuleIssue[] = [];

  const sortedOccurrences = [...event.occurrences].sort(
    (a, b) => new Date(a.startDate).getTime() - new Date(b.startDate).getTime(),
  );

  const pushIssues = (occurrence: (typeof sortedOccurrences)[number], ruleType: string, violations: IRuleViolation[]) => {
    for (const v of violations) {
      issues.push({
        eventId: event.id,
        eventName: event.name,
        occurrenceId: occurrence.id,
        occurrenceDescription: occurrence.description,
        academicYearStart,
        date: v.date,
        ruleType,
        fieldLabelKey: v.fieldLabelKey,
        messageKey: v.messageKey,
        messageParams: v.messageParams,
      });
    }
  };

  for (const occurrence of sortedOccurrences) {
    for (const builtInRule of [minDurationRule, minDaysToNextRule]) {
      pushIssues(occurrence, builtInRule.id, builtInRule.validate(event, occurrence, {}, context));
    }

    for (const rule of (event.rules ?? [])) {
      pushIssues(occurrence, rule.type, validateEventRule(rule, event, occurrence, context));
    }
  }

  return issues;
}
\end{lstlisting}

\clearpage
\section{Password Hashing and Verification} \label{app:crypto}

Referenced from Section~\ref{sec:backend-authentication}.

\begin{lstlisting}[caption={[Password hashing and verification, in full.]{Password hashing and verification, in full (\texttt{server/auth/crypto.ts}).}}, label={lst:crypto-full}]
export const hashPassword = (password: string): string => {
  const salt = crypto.randomBytes(16).toString("hex");
  const hash = crypto.scryptSync(password, salt, 64).toString("hex");
  return `scrypt$${salt}$${hash}`;
};

export const verifyPassword = (password: string, encodedHash: string): boolean => {
  const [algorithm, salt, storedHash] = encodedHash.split("$");
  if (algorithm !== "scrypt" || !salt || !storedHash) {
    return false;
  }

  const computedHash = crypto.scryptSync(password, salt, 64).toString("hex");
  if (computedHash.length !== storedHash.length) return false;
  return crypto.timingSafeEqual(
    Buffer.from(storedHash, "hex"),
    Buffer.from(computedHash, "hex"),
  );
};
\end{lstlisting}

\clearpage
\section{The withRolePolicy Authorization Wrapper} \label{app:with-role-policy}

Referenced from Section~\ref{sec:backend-rbac}.

\begin{lstlisting}[caption={[The withRolePolicy authorization wrapper, in full.]{The \texttt{withRolePolicy} authorization wrapper, in full (\texttt{server/shared/authorize.ts}).}}, label={lst:with-role-policy-full}]
export function withRolePolicy<
  TInstance extends object,
  TPolicy extends Partial<Record<keyof TInstance, TRolePolicy>>,
>(
  instance: TInstance,
  policy: TPolicy,
  getRole: (args: unknown[]) => TUserRole,
): { [K in keyof TPolicy]: TInstance[K & keyof TInstance] } {
  const wrapped: Record<string, unknown> = {};

  for (const key of Object.keys(policy)) {
    const minimum = policy[key as keyof TPolicy] as TRolePolicy;
    const original = instance[key as keyof TInstance] as unknown as
      (...args: unknown[]) => unknown;

    wrapped[key] = (...args: unknown[]) => {
      if (minimum !== "any") {
        assertRole(getRole(args), minimum);
      }
      return original.apply(instance, args);
    };
  }

  return wrapped as { [K in keyof TPolicy]: TInstance[K & keyof TInstance] };
}
\end{lstlisting}

\clearpage
\section{The DraggableEvent Component} \label{app:draggable-event}

Referenced from Section~\ref{sec:frontend-calendar-ui}.

\begin{lstlisting}[caption={[The DraggableEvent wrapper component, in full.]{The \texttt{DraggableEvent} wrapper component, in full (\texttt{features/calendar/dnd/draggable-event.tsx}).}}, label={lst:draggable-event-full}]
export function DraggableEvent({
  event,
  occurrence,
  children,
  className,
}: DraggableEventProps) {
  const { canEditEvents } = useCalendar();
  const { startDrag, endDrag, isDragging, draggedOccurrence } = useDragDrop();

  if (!canEditEvents) {
    return <>{children}</>;
  }

  const isCurrentlyDragged = isDragging && draggedOccurrence?.occurrence.id === occurrence.id;

  const handleClick = (e: React.MouseEvent<HTMLDivElement>) => {
    e.stopPropagation();
  };

  return (
    <motion.div
      className={`${className || ""} ${isCurrentlyDragged ? "opacity-50 cursor-grabbing" : "cursor-grab"}`}
      draggable
      onClick={(e: React.MouseEvent<HTMLDivElement>) => handleClick(e)}
      onDragStart={(e) => {
        (e as DragEvent).dataTransfer!.setData(
          "text/plain",
          occurrence.id,
        );
        startDrag(event, occurrence);
      }}
      onDragEnd={() => {
        endDrag();
      }}
    >
      {children}
    </motion.div>
  );
}
\end{lstlisting}
